\section{Introduction and theoretical grounding}
Parameter estimation for the TC-Lab must distinguish uncertainty in a measured
thermal response from ambiguity in its physical explanation. A model may predict
the observed temperature closely while several parameter combinations produce
the same response. Changing parameter coordinates can improve numerical scaling,
but an invertible coordinate transformation preserves that ambiguity.

Continuous-time identification also distinguishes the input-driven model from
the residual process. The latter can include measurement error, unmeasured
disturbances, and model inadequacy \citep{garnier-2008-direct}. A fitted residual
covariance therefore does not by itself identify sensor noise or establish a
physical disturbance mechanism. This distinction guides the uncertainty
comparisons below: we vary working noise assumptions and examine the errors
that remain on experiments excluded from calibration.

Identifiability concerns the uniqueness of parameters given a model, an observed
output, and an experiment. Practical estimability additionally depends on the
record length, input, noise, and initial-state information. Profile likelihood
examines a parameter while reoptimizing the remaining parameters and can reveal
features missed by a local quadratic approximation
\citep{raue-2009-structural}. We combine that approach with explicit transfer
functions, sensitivity singular values, and resampling. The algebraic calculation
is needed to distinguish an exact equivalence class from a numerically shallow
objective or a finite exploration bound.

Linear time-invariant system identification supplies a connection between input
excitation, parameter information, and prediction accuracy. Application-oriented
design can target informative parameter combinations rather than require equal
precision in every physical parameter. \citet{hjalmarsson-2009-system} develops
this distinction using parameter confidence regions and application loss. Its
prediction-error information formula assumes that the true system belongs to
the model set and that the predictor is differentiable. Those assumptions require
separate scrutiny here because the TC-Lab residuals can contain serial structure
and between-experiment discrepancies.

Finite-horizon amplitude constraints introduce a further distinction between an
information objective and its numerical solution. \citet{manchester-2010-input}
formulates time-domain information design as a nonconvex input problem and
develops a convex relaxation. We use bounded multistart nonlinear optimization
with explicit command and temperature constraints. We have not implemented that
relaxation, so its approximation guarantees do not apply to the designs reported
here. Comparing local solutions across assumed models and residual covariances
is part of assessing a proposed experiment.

An $n$-stage chain is a family of distributed-delay models. The linear chain
construction relates cascaded first-order stages to Erlang dwell times
\citep{hurtado-2019-generalizations}. For a fixed mean delay, increasing the stage
count narrows the delay distribution and approaches a pure delay. Fitting this
family therefore tests a shape assumption as well as a numerical approximation.
Its auxiliary stage count differs from its number of estimated parameters.

This study asks how these distinctions affect the available TC-Lab data. We
compare physical and input--output parameterizations, quantify conditional
uncertainty, test sensitivity to initialization and data handling, and evaluate
simulated inputs for future estimation and discrimination. The study is a
reproducible draft for internal review; it does not establish a validated
physical mechanism from temperature fit alone.
